\section{Arbitrary continuous obstacles on finite smooth quotient germs}
\label{sec:quotient}
\textit{Support: the entire qualified G331 and earlier own n1,n2. This is new scoped original-residual progress, not a theorem for every excellent normal germ.}

Suppose $S$ is a positive-dimensional smooth affine closed-germ local domain over an algebraically closed characteristic-zero field, and a finite group $G$ acts faithfully by local automorphisms. Assume $R=S^G$ is excellent Noetherian and $S$ is finite over $R$; these hypotheses hold for the algebraic quotient example below. The ring $R$ is a normal local domain: an element of $\operatorname{Frac}(R)$ integral over $R$ is integral over $S$, hence belongs to the normal ring $S$, and, being invariant, belongs to $R$. Fix a proper primary $\mathfrak a\subset R$, put $b=\mathfrak aS$, and let
\[
 \rho:V_b(S)\longrightarrow V_{\mathfrak a}(R)
\]
be restriction. This map is continuous surjective by G331, including semivaluations with nonminimal supports. Both the obstacle and the normalizing source ideal $b$ are invariant.

We need the entire fibres, not just field valuations. Let $K=\operatorname{Frac}(R)$ and $L=\operatorname{Frac}(S)$. Then $L^G=K$: if $s/t\in L^G$, multiply its denominator by the finite product of its $G$-translates; that product lies in $R$, and the resulting invariant numerator lies in $R$. Here is the needed fixed-field argument, so no unpublished fibre-transitivity theorem is assumed. Distinct automorphisms $\sigma_i$ are linearly independent as $L$-valued functions on $L$: in a relation with the fewest nonzero coefficients, evaluate at $ax$ and subtract $\sigma_j(a)$ times the relation at $x$. Choosing $a$ on which two distinct automorphisms differ produces a shorter nonzero relation, a contradiction; a one-term relation is impossible by evaluating at $1$. Consequently the evaluation columns $(\sigma_i(x))_i$ span $L^d$, for otherwise a nonzero annihilating row would give such a relation. Choose $x_1,\ldots,x_d$ with invertible evaluation matrix. For arbitrary $x\in L$, solve
\[
 \sigma_i(x)=\sum_j c_j\sigma_i(x_j)\quad\text{for every }\sigma_i\in G.
\]
Applying any $\tau\in G$ permutes the rows, so uniqueness forces $\tau(c_j)=c_j$. Hence $c_j\in K$, and the identity row writes $x=\sum_jc_jx_j$. The invertible matrix also proves the $x_j$ are $K$-linearly independent. Thus $[L:K]=d$. For each $x$, the product over its distinct orbit elements of $T-\sigma(x)$ is a separable polynomial with invariant coefficients. Its minimal polynomial over $K$ divides it and splits in $L$. This proves $L/K$ is Galois, with precisely $G$ as its automorphism group. Also $S$ is the integral closure of $R$ in $L$, since any element integral over $R$ is integral over $S$ and lies in $S$ by normality.

Use now the full qualified G331 normal-target coefficient-norm trace theorem, including its proof of full fibre support on all semivaluations. For any $w\in\rho^{-1}(v)$ one may take the splitting field $L$ and the lift $w$ itself on its finite integral closure $S$. There is no inseparable multiplicity in characteristic zero. That theorem says the canonical trace measure is
\[
 \mu_v=\sum_{\sigma\in G}\delta_{w\circ\sigma},\qquad
 \operatorname{supp}\mu_v=\rho^{-1}(v). \tag{9}
\]
Thus each restriction fibre is exactly the $G$-orbit of any one of its points, with repeated group points counted in the measure. This use of G331 is essential: its coefficient identity is proved using valued quotient fraction fields after killing supports. We do not apply Galois extension theory to an original fraction field through a nonzero semivaluation support. Formula (9) retains all those supports and proves the needed full-fibre orbit assertion without deleting a singular stratum.

For every finite continuous signed obstacle $f$ on $V_{\mathfrak a}(R)$ set $F=f\circ\rho$. By n2 its greatest source geometric minorant $H_S=P_S(F)$ belongs to the continuous class $C_S$. Pullback by each automorphism preserves $C_S,D_S$ and $F$. By uniqueness of the greatest minorant, $H_S\circ\sigma=H_S$. Formula (9) makes $H_S$ constant on every restriction fibre. G331's positive continuous descent therefore yields a unique $H_R\in C_R$ with $\rho^*H_R=H_S$.

Surjectivity gives $H_R\le f$. If $u\in D_R$ and $u\le f$, its positive pullback lies in $D_S$ below $F$, so $\rho^*u\le H_S=\rho^*H_R$. Surjectivity again gives $u\le H_R$. Consequently
\[
 P_R(f)=H_R\in C_R,\qquad \rho^*P_R(f)=P_S(f\circ\rho), \tag{10}
\]
for every such $f$. This is actual geometric positivity and continuity, not positivity ascribed to an upper semicontinuous regularization. It applies to arbitrary continuous obstacles, including obstacles that are not model functions and without any canonical finite-generation hypothesis.

There is an explicit sequence of downstairs actual primary Cartier profiles. Let $B_m=A_m^S(F)$. These ideals are $G$-stable: replacing $w$ by $w\circ\sigma$ in their defining all-support inequalities proves this directly. Their profiles $h_m^S=-w(B_m)/m$ are therefore orbit constant. Define $J_m=\mathcal N_{L/K}(B_m)$, the coefficient norm ideal of G331. To spell out the construction, choose finite $R$-module generators $s_i$ of $B_m$, form
\[
 N_{L(T)/K(T)}\Big(\sum_i T_i s_i\Big)
 =\prod_{\sigma\in G}\Big(\sum_i T_i\sigma(s_i)\Big),
\]
and let $J_m$ be its coefficient ideal in $R$. Coefficients are integral over normal $R$ and lie in $K$, hence lie in $R$. G331 proves independence of generators and the all-support identity
\[
 v(J_m)=\sum_{\sigma\in G}(w\circ\sigma)(B_m)=d\,w(B_m). \tag{11}
\]
It also proves the ideal contains a primary power: $B_m$ contains $\mathfrak a^NS$, and pure-variable coefficients for base generators contain their $d$th powers. If the envelope is nonzero, $B_m$ is proper and (11) makes $J_m$ proper primary; the unit case gives the zero profile. Hence
\[
 \kappa_m=-v(J_m)/(dm),\qquad \rho^*\kappa_m=h_m^S,
\]
are actual geometric primary profiles below $f$ and converge uniformly to $P_R(f)$. This follows from n2's smooth continuity and n1's canonical uniform criterion; the source-to-target restriction is surjective, so the sup norms agree. The construction involves no finite flatness or determinant norm on an arbitrary finite module. It uses only the granted normal-target coefficient norm. The canonical ideals downstairs are also exactly $A_m^R(f)=B_m\cap R$, by the same all-support restriction argument as n1; we do not infer finite generation from this contraction.

For every signed real Cartier model obstacle on an actual primary model over $R$, a stronger bound is available. G331's exact evaluation formula for an integral exceptional Cartier divisor expresses it as a finite signed sum of primary ideal values with a fixed primary section correction; real finite presentations give real finite sums. Pulling these ideals to $S$ and using the smooth common principalization makes the maximal-ideal-gauge source obstacle a function determined on one BFJ dual complex. In n2's notation let $M,C$ be its fixed model constants and set
\[
 e=CM/\min_{\nu\in V_{\mathfrak m_S}(S)}\nu(b).
\]
The denominator is positive and bounded by primaryness. The exact gauge transport turns (8) into
\[
 0\le P_R(f)-\kappa_m\le e/m \quad\text{on all }V_{\mathfrak a}(R). \tag{12}
\]
In particular $J_m\subset A_{dm}^R(f)$ and $h_{dm}^R\ge\kappa_m$. To bound the \emph{full} downstairs canonical sequence, for $k\ge d$ put $m=\lfloor k/d\rfloor$ and $r=dm$, so $0\le k-r<d$. Grading with $(A_1^R)^{k-r}$ and nonpositivity gives
\[
 h_k^R\ge h_r^R-dM_1^R/k.
\]
Combining with (12), and using $m\ge k/(2d)$ for $k\ge2d$, gives
\[
 0\le P_R(f)-h_k^R
 \le \frac{e}{\lfloor k/d\rfloor}+\frac{dM_1^R}{k}
 \le\frac{2de+dM_1^R}{k}\qquad(k\ge2d). \tag{13}
\]
Every ideal used in (12)--(13) is unit or genuine primary in every degree. The constants depend on this obstacle and smooth model; the graded algebra need not be finitely generated. We assert no new finite-generation example or general excellent-germ bound. The primary correction is fixed for the obstacle, and all infinite generator values are accommodated by its finite primary-ideal values.

Here is an actual normal nonisolated higher-dimensional domain in the new theorem. For $r\ge1$, take
\[
 S=k[x,y,z_1,\ldots,z_r]_{(x,y,z_1,\ldots,z_r)},\quad
 G=\{1,\sigma\},\quad \sigma(x,y,z)=(-x,-y,z).
\]
The invariant monomials have even total $x,y$ degree and are generated by $a=x^2,b=xy,c=y^2,z_i$. Thus
\[
 R=k[a,b,c,z_1,\ldots,z_r]_{(a,b,c,z_1,\ldots,z_r)}/(ac-b^2).
\]
The only relation is $ac=b^2$: monomials reduced by this relation have distinct images, and the image has the displayed dimension. The extension is finite, since $x,y$ satisfy their quadratic equations; localizing the invariant ring at the origin gives exactly the indicated local cover. To check localization of invariants explicitly, multiply any local denominator by its orbit product; the product is an invariant unit at the origin, and an invariant fraction then has invariant numerator. The cover's unique point above the origin is the origin. The rings are excellent finite-type local rings. Normality follows from the invariant argument already given. The dimension is $r+2$, and the singular locus is $(a,b,c)$, of dimension $r$: off it a derivative of $ac-b^2$ is a unit, whereas on it the local quadratic equation has no linear term and embedding dimension exceeds dimension by one. This is a normal nonisolated germ, not a nonnormal crossing example.

For a concrete continuous signed obstacle with cancellation sensitivity take $r=1$, $\mathfrak m=(a,b,c,z)$,
\[
 I=(a+b,b+c,z)+\mathfrak m^2,\qquad
 J=(a,c,z^2)+\mathfrak m^3,\qquad
 f(v)=v(I)-5/4+\tfrac1{100}\sin(v(J)). \tag{14}
\]
Both ideals are proper primary, so (14) is continuous and finite on every support. On the smooth cover substitute $x=r_0$, $y=r_0(h-1)$, retaining $z$, and use monomial weights $(1,1,6)$ on $(r_0,h,z)$. This defines a valuation after the injective substitution; its restriction gives $\mathfrak m$ value two. The sums $a+b=x(x+y)$ and $b+c=y(x+y)$ have weight three, so after normalization $v(I)=3/2$, while $v(J)=1$. A monomial valuation with $(x,y,z)$ weights $(1,1,6)$ gives the same values on $a,b,c,z$ but $v(I)=1$, $v(J)=1$. Thus the respective obstacles are $1/4+\sin(1)/100>0$ and $-1/4+\sin(1)/100<0$. They detect cancellation despite identical coordinate weights. The envelope of (14) is constructed by the $J_m$ profiles above and is continuous geometric; no closed formula or claim about finite generation for this example is made. Formula (10) applies to every other continuous signed obstacle on this germ as well, independently of any finite contact certificate.

Equations (10)--(13) obey the exact primary-gauge cocycle of n1; the canonical ideals are unchanged by gauge transport. The quotient map is not a normalization map of the singular invariant ring: that ring is already normal. Its fibres need not be singleton, and their orbit compatibility is proved rather than assumed. For arbitrary nonnormal germs only n1's full-fibre conditional normalization rule is obtained. The excellent-curve comparison, nilpotent reduction, empty dimension-zero case, and literal radial distinction remain as in n1. This section constructs no mixed Monge--Ampere product, energy, or unrestricted finite-map operator.
